\section{Future Work}
\label{sec:concl-future}

The immediate priority is to separate secondary magnetic loss from
input-path heating through matched current, temperature, voltage-drop,
and power measurements. The suspected single-via restriction should be
verified before additional copper or diode relocation is selected.
Turn-on voltage should be measured together with gate timing and resonant current.

The approximately \SI{85}{\watt} V2 result needs to be repeated using the
original hardware and settings \cite{projectOctoberPlan}.
The \SI{79.8928}{\watt} result belongs to the older assembly.
Synchronized input, secondary-bus, connector, and auxiliary power,
with component temperature histories, would locate the main losses.
Long-duration tests should start at the \SI{65}{\watt} minimum under
defined electrical and thermal conditions.

The diagnostic tests in Section~\ref{sec:disc-plan-gap} should guide
individual changes and matched retests toward \SI{100}{\watt}
continuous. Any limiting result should report repeatable output, duration,
and measured stop condition. Attachment, detachment, overload, and fault
tests require synchronized secondary voltage, VBUS, board-specific enable,
and firmware state against declared bands and recovery times.

An independent secondary shutdown or discharge path should be developed
alongside a model of installed magnetics, switching parasitics, packet
energy, and commutation. This would support evaluation of scheduling,
capacitance, layout, and cooling changes.

V3 requires completion around the selected magnetics and external output
stage, followed by review, assembly, and gradual bring-up. Comparison must
use equivalent functions, volume definitions, supported power, cooling,
and electrical conditions.

USB-PD fixed-PDO and PPS protocol tests, independent safety assessment,
and EMC verification are also required before assigning a supply rating.
